\PassOptionsToPackage{table,xcdraw}{xcolor}
\documentclass[11pt, a4paper]{article}

\usepackage[T1]{fontenc}
\usepackage[utf8]{inputenc}
\usepackage{tgpagella}
\usepackage[spanish, es-noshorthands]{babel}
\usepackage{amsmath, amssymb}
\usepackage[most]{tcolorbox}
\usepackage[margin=2.5cm]{geometry}
\usepackage{fancyhdr}

\pagestyle{fancy}
\fancyhf{}
\setlength{\headheight}{16pt}
\lhead{\textbf{UCB Sede Tarija} | Ing. Mecatrónica}
\rhead{IMT-221 | Requerimiento Pre-Lab Hito 2}
\renewcommand{\headrulewidth}{0.4pt}
\definecolor{ucbblue}{RGB}{0, 51, 102}

\begin{document}

\begin{tcolorbox}[colback=ucbblue!5, colframe=ucbblue, title=\textbf{Requisito de Avance: Entre Sesiones (Pre-Lab Hito 2)}]
    Para garantizar que la implementación del hardware y las derivaciones matemáticas están alineadas, cada grupo debe remitir el siguiente informe de progreso antes de la próxima sesión[cite: 18].
\end{tcolorbox}

\section*{Lista de Entregables Exigidos}
Envíe un documento PDF respondiendo de forma explícita a los siguientes 10 puntos[cite: 18]. Si su grupo no dispone de hardware físico (Track B), complete los campos teóricos/simulados y marque los medidos como explícitamente "Pendiente"[cite: 18].

\begin{enumerate}
    \item \textbf{Evidencia Fotográfica / Esquemática:} Fotografía nítida del estado actual del circuito en protoboard, o el esquema desarrollado detalladamente[cite: 18].
    \item \textbf{Corriente de Cola ($I_T$):} Valor medido físicamente de la corriente generada por el espejo Q3-Q4 (si el hardware está disponible)[cite: 18].
    \item \textbf{Estado Balanceado DC:} Mediciones físicas logradas en el Par Diferencial (Voltajes $V_E$, $V_{C1}$, $V_{C2}$) asumiendo entradas $0\,\text{V}$[cite: 18].
    \item \textbf{Evidencia de Simulación (LTspice):} Captura del esquemático simulado demostrando el Punto de Operación (.op) del circuito completo[cite: 18].
    \item \textbf{Identidad del Transistor:} Registro exacto del modelo/número de parte (ej. 2N3904) del transistor empleado tanto en hardware como en la simulación LTspice (proveer la directiva \texttt{.model} utilizada)[cite: 18].
    \item \textbf{Transconductancia ($g_m$):} Cálculo matemático explícito de $g_m$ utilizando la corriente $I_C$ de reposo real/simulada de su circuito (No la teórica ideal)[cite: 18].
    \item \textbf{Ganancia Diferencial ($A_d$):} Predicción de su ganancia $A_d \approx g_m R_C$ utilizando los valores reales de $R_C$ medidos por su grupo[cite: 18].
    \item \textbf{Problemas no resueltos:} Descripción de cualquier comportamiento anómalo detectado (ej. offset severo, saturación)[cite: 18].
    \item \textbf{Componentes Faltantes:} Lista de cualquier elemento del BOM que el grupo aún deba adquirir[cite: 18].
    \item \textbf{Próxima Acción Física:} Describa cuál será la siguiente conexión o medida a realizar en el laboratorio[cite: 18].
\end{enumerate}

\end{document}
